\thispagestyle{plain}
\noindent
\begin{tikzpicture}
\node[fill=fnDark,text=white,rounded corners=2pt,inner sep=6pt,
      minimum width=\textwidth,anchor=west]
  {\large\bfseries FPGA-Neural V2 --- General description and features};
\end{tikzpicture}

\vspace{6pt}
\noindent
{\small FPGA-Neural V2 is a \textbf{neural multiprocessor / dataflow machine},
the evolution of the V1 sequential accelerator (documented separately, frozen
and unmodified as the project's golden reference). Where V1 executes one
neuron at a time under host-driven SPI control, V2 registers a
\textbf{dependency graph of neurons} and keeps \code{N\_SLOTS} independent
Neural Processor $+$ Memory Manager pairs busy concurrently, resolving data
dependencies and hiding PSRAM latency in hardware, without host
intervention once a graph is loaded. Computation (INT8 MAC, ReLU,
saturation) is bit-exact identical to V1's own datapath; what changed is
everything \emph{around} it.}

\vspace{8pt}
\begin{multicols}{2}
{\color{fnDark}\large\bfseries Features}\\[2pt]
{\footnotesize
\begin{itemize}[leftmargin=1.1em]
\item \textbf{Dependency-graph scheduling}: nodes are registered with an
      explicit producer list; a node becomes eligible for execution only
      once every producer it depends on has genuinely completed --- verified
      for 1-hop shared-producer/multi-consumer graphs and 2-hop transitive
      (diamond) graphs.
\item \code{N\_SLOTS} independent \textbf{Neural Processor + Memory Manager}
      pairs (default recommended: \textbf{2}), each running the identical
      8-stage INT8 pipeline inherited from V1.
\item \textbf{Word-level burst memory backend}: fetches move a full 16-bit
      PSRAM word per transaction instead of one byte, reusing
      \code{memory\_interface.v}/\code{psram\_controller.v} directly and its
      already-implemented page-mode support --- \textbf{2.24--2.37$\times$}
      real wall-clock speedup, measured.
\item \textbf{Shared on-chip activation cache}: a vector of activations
      shared by many neurons of the same layer is fetched from PSRAM
      \emph{once}, not once per neuron --- a further real
      \textbf{1.66--2.00$\times$} cycle reduction on shared-input workloads.
\item Same \textbf{real, unmodified V1 PSRAM backend} throughout
      (\code{memory\_interface.v}, \code{psram\_controller.v}) --- V1 remains
      the frozen golden reference and was never altered to make V2 look
      faster.
\item \textbf{Real, measured} characterization at every step: Verilator
      RTL simulation, Yosys synthesis, real \code{nextpnr-ecp5}
      place\&route --- no theoretical number reported without a matching
      real measurement.
\end{itemize}}

\columnbreak

{\color{fnDark}\large\bfseries Honest, measured limitations}\\[2pt]
{\footnotesize
\begin{itemize}[leftmargin=1.1em]
\item The system is \textbf{memory-bound}, not compute-bound: real compute-
      to-memory-wait ratio on the order of 1:170--1:220. A single shared
      PSRAM port saturates at $\approx$90\% utilization regardless of
      \code{N\_SLOTS}$\ge$2 --- real parallel scaling beyond 2 slots is
      essentially flat for large workloads.
\item \code{N\_SLOTS=4} is \textbf{not recommended}: it delivers no
      additional real throughput once the shared PSRAM port saturates,
      and with the activation cache active it \textbf{fails the 80\,MHz
      timing target outright} (65.01\,MHz measured).
\item Fixed, lowest-index-priority arbitration (Director and memory
      arbiter alike) is not fairness-balanced --- a real, measured
      per-slot workload imbalance exists under sustained contention.
\end{itemize}}

\vspace{4pt}
{\color{fnDark}\large\bfseries Target \& toolchain}\\[2pt]
{\footnotesize
\begin{itemize}[leftmargin=1.1em]
\item FPGA: Lattice ECP5 \code{LFE5U-45F-8BG381C} ($-8$, CABGA381) --- same
      target device as V1.
\item Synthesis: Yosys; place\&route: real \code{nextpnr-ecp5}.
\item Simulation: Verilator 5.050 (\code{--binary --timing}) --- adopted
      for V2 after two independent Icarus Verilog v13.0 scheduling
      defects were found and reproduced on minimal repros (V1's own
      certification, performed separately, was unaffected).
\item PSRAM: ISSI \code{IS66WVE4M16EBLL-70BLI} (64\,Mb, 4M$\times$16),
      real chain reused byte-for-byte from V1.
\end{itemize}}
\end{multicols}

\vspace{2pt}
% --- key parameter table ---
\noindent
{\small\color{fnDark}\bfseries Key parameters (recommended configuration, real measured data)}
\vspace{2pt}

\noindent
\begin{tabularx}{\textwidth}{L{3.6cm}L{3.6cm}Y}
\toprule
\rowh \thd{Quantity} & \thd{Value} & \thd{Notes} \\
\midrule
Data precision & INT8 (signed) & \code{DATA\_WIDTH}=8, identical to V1 \\
\rowa Accumulator & INT32 (signed) & \code{ACC\_WIDTH}=32 \\
Dot-product width & 8 & \code{P\_IN}=8 parallel MAC lanes per neuron \\
\rowa Recommended concurrency & \code{N\_SLOTS}=2 & real, measured net win; see ch.~\ref{ch:impl2} \\
Fmax, full system (\code{N\_SLOTS}=2) & 87.72~MHz & real place\&route, word-burst + activation cache active \\
\rowa cycles/neuron (1 neuron, 8 inputs, real PSRAM) & 166 (V1: 209) & \textbf{2.6$\times$} real wall-clock speedup vs V1 \\
Combined real speedup vs baseline (\code{N\_SLOTS}=2) & \textbf{2.45$\times$} & word-burst $+$ activation cache, D-Stress workload \\
\rowa Address space & 23~bit (byte) & \code{ADDR\_WIDTH}=23, unchanged from V1 \\
\bottomrule
\end{tabularx}

\vspace{8pt}
\noindent
{\small\color{fnDark}\bfseries System block diagram}
\begin{center}
\begin{tikzpicture}[node distance=6mm and 9mm,font=\footnotesize]
  \node[fnblockD,minimum width=24mm,minimum height=15mm] (host){HOST\\{\scriptsize registers a node graph}};
  \node[fnblockT,right=14mm of host,minimum width=30mm,minimum height=13mm] (dm){Dependency\\Manager};
  \node[fnblockT,right=14mm of dm,minimum width=28mm,minimum height=13mm] (dir){Neural\\Director};
  \node[fnreg,fill=white,right=14mm of dir,minimum width=30mm,minimum height=20mm] (slots){
    \begin{tabular}{c}
      N\_SLOTS $\times$ \\
      Memory Manager \\
      $+$ Neural Processor
    \end{tabular}};
  \node[fnblockA,below=9mm of dir,minimum width=28mm,minimum height=11mm] (cache){Activation\\Cache};
  \node[fnblock,right=14mm of slots,minimum width=22mm,minimum height=15mm] (ram){PSRAM 8\,MB\\{\scriptsize real V1 backend}};
  \draw[fnbus] (host) -- (dm);
  \draw[fnbus] (dm) -- node[fnlbl,above]{ready node} (dir);
  \draw[fnbus] (dir) -- (slots);
  \draw[fnarrowT] (slots.south) |- (cache.east);
  \draw[fnarrowT] (cache.north) |- node[fnlbl,above]{producer done} (dm.south);
  \draw[fnbus] (slots) -- node[fnlbl,above]{16-bit word} (ram);
  \draw[fnbus] (cache.east) -- ++(6mm,0) |- ([yshift=-2mm]ram.south);
\end{tikzpicture}
\end{center}
\begin{center}\footnotesize\itshape\color{fnGrey}
A slot's completion feeds back to the Director (frees the slot) and to the
Dependency Manager (wakes up any node waiting on it) --- closing the
dataflow loop entirely on-chip.\end{center}
